\chapter{Dengue Fever Antibody Test}

\section*{What Is This Test?}
This test detects infection by the Dengue virus, a flavivirus transmitted by \textit{Aedes} mosquitoes. It measures viral antigen and specific antibodies to diagnose dengue fever in patients presenting with acute febrile illness \cite{gubler1998dengue}.

\section*{Alternative Names}
Dengue serology, Dengue IgM and IgG, Dengue NS1 antigen test, Dengue PCR

\section*{Principle}
Diagnostic markers vary by the phase of illness. During the acute phase (days 1 to 7), the viral non-structural protein 1 (NS1) antigen is highly detectable in blood using ELISA. Beyond day 4, IgM antibodies rise, indicating acute or recent infection. IgG antibodies develop later and persist long-term, indicating past infection or secondary dengue.

\section*{History}
Dengue-like outbreaks were recorded as early as 992 AD in China. The virus was isolated in 1943 during World War II \cite{guzman2015dengue}. Serological assays using ELISA to detect NS1 antigen and IgM/IgG antibodies became widely available in the early 2000s, revolutionizing rapid diagnosis in endemic regions.

\section*{Why This Test Is Ordered}
This test is ordered if you present with sudden high fever, severe headache, retro-orbital (behind-the-eye) pain, severe joint and muscle pain ("breakbone fever"), rash, or minor bleeding, particularly after traveling to tropical or subtropical regions where dengue is endemic \cite{who2009dengue}.

\section*{Primary vs Secondary Test}
This is a primary test for diagnosing acute or past Dengue virus infection.

\section*{How This Test Is Performed}
A blood sample is drawn from a vein in your arm.

\section*{What Preparation Is Needed}
No fasting or special preparation is required.

\section*{Contraindications and Precautions}
\begin{itemize}
  \item Standard venipuncture precautions apply. Interpretation must account for cross-reactivity with other flaviviruses (such as Zika virus, West Nile virus, or Yellow Fever vaccine).
\end{itemize}
\section*{Risks}
Minimal risk associated with standard blood collection (minor pain, bruising).

\section*{What Happens After the Test}
Dengue fever has no specific antiviral treatment. Treatment is supportive (hydration, rest). If the test is positive, avoid aspirin and NSAIDs (like ibuprofen) because they increase the risk of severe bleeding complications (dengue hemorrhagic fever); use acetaminophen for pain and fever control instead.

\section*{Understanding Results}

\subsection*{Below Normal}
A negative result (no NS1 antigen or antibodies detected) suggests that Dengue virus is not the cause of your illness. However, if the blood was drawn too early, antibodies may not yet have developed.

\subsection*{Above Normal}
- **Positive NS1 or IgM:** Confirms an active or very recent Dengue virus infection.
- **Positive IgG only:** Indicates a past Dengue virus infection.
- **Highly elevated IgG with positive IgM:** Suggests a secondary dengue infection, which carries a higher risk of severe dengue complications \cite{simmons2012dengue}.

\section*{Normal Results}
Negative.

\section*{Follow-up Tests}
A complete blood count (CBC) should be performed regularly to monitor for thrombocytopenia (low platelet count) and hemoconcentration (rising hematocrit), which are key warning signs of progression to severe dengue \cite{cdc2023dengue}.

\section*{References}
\bibliographystyle{unsrtnat}
\bibliography{references}

\clearpage
\section*{Regulatory Status and Authorization Date}
This chapter describes a test category, not a specific commercial product. FDA, CDSCO, EU IVDR, and MHRA approval or registration dates therefore cannot be assigned without the assay manufacturer, product name, instrument, and intended use. For laboratory-developed tests, an individual agency approval date may not exist. See the Preface for the interpretation of regulatory status.

\clearpage
